\section{Valuations}

\begin{definition}[Valuation]
    Let \(K\) be a field. A \emph{valuation} on \(K\) is a function
    \[
        v:K^\times\longrightarrow \mathbb R,
    \]
    where \(K^\times=K\setminus\{0\}\), satisfying the following conditions:
    \begin{enumerate}
        \item For all \(a,b\in K^\times\),
              \[
                  v(ab)=v(a)+v(b).
              \]
        \item For all \(a,b\in K^\times\) with \(a+b\neq 0\),
              \[
                  v(a+b)\geq \min\{v(a),v(b)\}.
              \]
    \end{enumerate}
    We extend \(v\) to \(K\) by setting
    \[
        v(0)=+\infty.
    \]
    With this convention, the second condition can be written as
    \[
        v(a+b)\geq \min\{v(a),v(b)\}
    \]
    for all \(a,b\in K\).
\end{definition}

\begin{proposition}
    Let \(v\) be a valuation on a field \(K\). Then
    \[
        v:(K^\times,\cdot)\longrightarrow (\mathbb R,+)
    \]
    is a group homomorphism.
\end{proposition}

\begin{proof}
    By definition, for all \(a,b\in K^\times\),
    \[
        v(ab)=v(a)+v(b).
    \]
    Also,
    \[
        0=v(1)=v(aa^{-1})=v(a)+v(a^{-1}),
    \]
    so
    \[
        v(a^{-1})=-v(a).
    \]
    Therefore \(v\) is a group homomorphism from \(K^\times\) to the additive group
    \((\mathbb R,+)\).
\end{proof}

\begin{definition}[Value group]
    Let \((K,v)\) be a valued field. The subgroup
    \[
        v(K^\times)\subseteq (\mathbb R,+)
    \]
    is called the \emph{value group} of \(v\).

    A field \(K\) together with a valuation \(v\) is called a \emph{valued field},
    and is denoted by
    \[
        (K,v).
    \]
\end{definition}

\begin{proposition}[Basic properties of valuations]
    Let \(v\) be a valuation on a field \(K\). Then for all \(a,b\in K^\times\), the
    following hold:
    \begin{enumerate}
        \item
              \[
                  v(1)=0,\qquad v(-1)=0.
              \]
        \item
              \[
                  v(-a)=v(a).
              \]
        \item
              \[
                  v(a^{-1})=-v(a).
              \]
        \item If
              \[
                  v(b)<v(a),
              \]
              then
              \[
                  v(a+b)=v(b).
              \]
    \end{enumerate}
\end{proposition}

\begin{proof}
    First,
    \[
        v(1)=v(1\cdot 1)=v(1)+v(1),
    \]
    so
    \[
        v(1)=0.
    \]
    Since \((-1)^2=1\), we have
    \[
        2v(-1)=v((-1)^2)=v(1)=0,
    \]
    and hence
    \[
        v(-1)=0.
    \]

    For \(a\in K^\times\),
    \[
        v(-a)=v((-1)a)=v(-1)+v(a)=v(a).
    \]

    Also,
    \[
        0=v(1)=v(aa^{-1})=v(a)+v(a^{-1}),
    \]
    so
    \[
        v(a^{-1})=-v(a).
    \]

    Finally, suppose that \(v(b)<v(a)\). By the valuation inequality,
    \[
        v(a+b)\geq \min\{v(a),v(b)\}=v(b).
    \]
    We claim that equality holds. If instead
    \[
        v(a+b)>v(b),
    \]
    then using
    \[
        b=(a+b)-a,
    \]
    we get
    \[
        v(b)\geq \min\{v(a+b),v(a)\}.
    \]
    But both \(v(a+b)\) and \(v(a)\) are strictly larger than \(v(b)\), so the right-hand
    side is strictly larger than \(v(b)\), a contradiction. Therefore
    \[
        v(a+b)=v(b).
    \]
\end{proof}




